
% ════════════════════════════════════════
% 符号说明（04.符号说明.tex）
% ════════════════════════════════════════
\section{符号说明}

主要符号及其含义如表\ref{tab:符号说明}所示。

\begin{table}[H]
  \centering
  \caption{符号说明}
  \label{tab:符号说明}
  {\small
  \renewcommand{\arraystretch}{0.95}
  \begin{tabular}{>{\centering\arraybackslash}p{0.30\textwidth}>{\centering\arraybackslash}p{0.10\textwidth}>{\centering\arraybackslash}p{0.52\textwidth}}
    \toprule
    符号 & 单位 & 说明 \\
    \midrule
    $k,\ \Delta t$ & —、h & 日内时间间隔序号（$k=0,1,\dots,143$）与时间步长（$1/6$ 小时，即 10 分钟） \\
    $L_k,\ S_k,\ S^{\tau}_k$ & kW & 第 $k$ 个间隔的小区负载与光伏实际功率，$S^{\tau}_k$ 为预报时刻 $\tau$ 发布的光伏预报功率 \\
    $\hat L_{d,k},\ \hat S_{d,k}$ & kW & 0:00 预测的该日第 $k$ 个间隔负载、光伏功率 \\
    $\varepsilon^{L}_{d,k},\ \varepsilon^{S}_{d,k}$ & kW & 负载、光伏的预测残差（实际值减预测值） \\
    $s,\ \omega_s,\ \pi_s$ & — & 情景序号、负载与光伏乘积因子对及情景概率（$s=1,\dots,9$） \\
    $L^{s}_{d,k},\ S^{s}_{d,k}$ & kW & 第 $s$ 个情景下第 $k$ 个间隔的负载、光伏功率 \\
    $p_k,\ g_k,\ g^{\mathrm{p}}_k,\ g^{\mathrm{a}}_k$ & 元/kWh、kW & 第 $k$ 个间隔的外网电价，及 0:00 计划购电功率与调整后购电功率 \\
    $c_k,\ d_k$ & kW & 第 $k$ 个间隔的储能充、放电功率 \\
    $E_k,\ E_{\min},\ E_{\max}$ & kWh & 第 $k$ 个间隔起始时刻的储电量及其上下限（1200、10800） \\
    $P_{\max},\ \eta$ & kW、— & 储能最大充放电功率（5000）与充放电效率（0.9） \\
    $\delta_k$ & kWh & 第 $k$ 个间隔的实际供电缺口（紧急购电能量） \\
    \bottomrule
  \end{tabular}}
\end{table}
